\documentclass[a4paper,12pt]{article}
\usepackage[utf8]{inputenc}
\usepackage[russian]{babel}
\usepackage{amsmath, amssymb}
\usepackage{geometry}
\geometry{top=2cm, bottom=2cm, left=2.5cm, right=1.5cm}

\begin{document}

\begin{center}

    {\normalsize \textit{ Федеральное государственное автономное образовательное }\\}
    {\normalsize \textit{учреждение высшего образования}\\}
    {\normalsize \textit{ «Национальный исследовательский университет ИТМО»}}

    \vspace*{4cm}

    {\Large \textbf{Практическая работа №3}}\\[0.3cm]
    {\large \textbf{Математическая статистика}}\\[1cm]

    \textit{Проверка гипотез. Параметрические критерии\\[3cm]}

    {\large
    Выполнили студенты: \\[0.1cm]
    Сафин Максим 467371 \\[0.1cm]
    Фокин Владимир 467883 \\[0.8cm]
    Поток: 22.3\\[0.8cm]
    Вариант: А-4
    }

    \vfill

    {\normalsize Санкт-Петербург, 2026}

\end{center}
\thispagestyle{empty}
\newpage

\section*{Постановка задачи}
Даны две независимые выборки из нормально распределённых генеральных совокупностей. Дисперсии неизвестны, но предполагаются равными. Проверяется гипотеза о равенстве математических ожиданий:
\[
H_0 : \mu_1 = \mu_2 \quad \text{против} \quad H_1 : \mu_1 \neq \mu_2.
\]
Объёмы выборок: $m = 54$, $n = 58$. Уровень значимости $\alpha = 0,05$. Используется двусторонний двухвыборочный $t$-критерий Стьюдента с объединённой оценкой дисперсии.

\section*{Исходные данные}
\textbf{Выборка X} (объём $m = 54$):
\[
\begin{array}{cccccccccc}
41 & 44 & 32 & 41 & 38 & 42 & 45 & 39 & 40 & 38 \\
38 & 33 & 41 & 45 & 48 & 41 & 55 & 26 & 42 & 42 \\
42 & 51 & 46 & 30 & 43 & 47 & 36 & 40 & 35 & 40 \\
38 & 44 & 31 & 46 & 37 & 40 & 39 & 43 & 42 & 46 \\
34 & 43 & 44 & 36 & 42 & 40 & 49 & 48 & 42 & 41 \\
38 & 45 & 43 & 35
\end{array}
\]

\textbf{Выборка Y} (объём $n = 58$):
\[
\begin{array}{cccccccccc}
45 & 46 & 50 & 38 & 41 & 43 & 41 & 49 & 45 & 47 \\
57 & 41 & 45 & 46 & 50 & 38 & 44 & 37 & 49 & 43 \\
47 & 45 & 34 & 47 & 52 & 50 & 44 & 40 & 50 & 47 \\
42 & 43 & 55 & 39 & 45 & 31 & 39 & 48 & 36 & 41 \\
45 & 39 & 34 & 51 & 42 & 44 & 38 & 51 & 43 & 40 \\
37 & 34 & 40 & 51 & 43 & 47 & 42 & 36
\end{array}
\]

\section*{Вычисление выборочных характеристик}

\subsection*{Выборка X}
Сумма значений:
\[
\sum_{i=1}^{42} x_i = 2207.
\]
Среднее арифметическое:
\[
\bar{x} = \frac{1}{m} \sum_{i=1}^{m} x_i = \frac{741}{54} \approx 40,8704.
\]
Сумма квадратов:
\[
\sum_{i=1}^{42} x_i^2 = 91703.
\]
Несмещённая выборочная дисперсия:
\[
s_1^2 = \frac{1}{m-1} \left( \sum x_i^2 - \frac{(\sum x_i)^2}{m} \right) = \frac{1}{54} \left( 91703 - \frac{2207^2}{54} \right) \approx 28,3414.
\]

\subsection*{Выборка Y}
Сумма значений:
\[
\sum_{j=1}^{45} y_j = 2527.
\]
Среднее арифметическое:
\[
\bar{y} = \frac{1}{n} \sum_{j=1}^{n} y_j = \frac{2527}{58} \approx 43,5690.
\]
Сумма квадратов:
\[
\sum_{j=1}^{58} y_j^2 = 111837.
\]
Несмещённая выборочная дисперсия:
\[
s_2^2 = \frac{1}{n-1} \left( \sum y_j^2 - \frac{(\sum y_j)^2}{n} \right) = \frac{1}{58} \left( 111837 - \frac{2527^2}{58} \right) \approx 30,4952.
\]

\vspace{0.1cm}
\section*{Объединённая оценка дисперсии}

Число степеней свободы:
\[
k = m + n - 2 = 54 + 58 - 2 = 110.
\]
Объединённая дисперсия:
\[
s^2 = \frac{(m-1)s_1^2 + (n-1)s_2^2}{k} = \frac{53 \cdot 28,3414 + 57 \cdot 30,4952}{110} \approx 29,4575.
\]

\section*{Стандартная ошибка разности средних}
\[
SE = \sqrt{s^2\left(\frac{1}{m} + \frac{1}{n}\right)} = \sqrt{29,4575 \cdot \left(\frac{1}{54} + \frac{1}{58}\right)} \approx \sqrt{1,0534} \approx 1,02565.
\]

\section*{Наблюдаемое значение t-статистики}
\[
t_{\text{набл}} = \frac{\bar{x} - \bar{y}}{SE} = \frac{40,8704 - 43,5690}{1,02565} \approx -2,6311.
\]
Абсолютное значение: $|t_{\text{набл}}| \approx 2,6311$.

\section*{Критическая область и принятие решения}
Для двустороннего критерия с уровнем значимости $\alpha = 0,05$ и числом степеней свободы $k = 110$ находим квантиль распределения Стьюдента:
\[
t_{1-\alpha/2;\, k} = t_{0.975;\, 110} \approx 1.9818.
\]
Критическая область:
\[
W = (-\infty; -1.9818] \cup [1.9818; +\infty).
\]
Так как $|t_{\text{набл}}| = 2,6311 > 1.9818$, наблюдаемое значение попадает в критическую область.

\section*{Вывод}
Мы отвергаем $H_0$, так как $|t_{\text{набл}}| = 2,6311 > t_{\text{крит}} = 1,982$. Теоретически $H_0$ может быть верна, но вероятность этого при полученных данных меньше 5\% (точнее, около 0,96\%). Поэтому мы считаем различие статистически значимым..

\end{document}